% ==================================================
% KATA PENGANTAR
% ==================================================

\chapter*{KATA PENGANTAR}

\phantomsection
\addcontentsline{toc}{chapter}{KATA PENGANTAR}

\normalfont
\justifying

Puji syukur penulis panjatkan kepada Allah SWT atas rahmat dan karunia-Nya
sehingga tesis ini dapat diselesaikan sebagai salah satu syarat untuk
menyelesaikan Program Magister Informatika di Institut Teknologi Bandung.

Penulis menyampaikan terima kasih kepada:

\begin{enumerate}
    \item Bapak Sugimin dan Ibu Umilah, selaku orang tua penulis, atas doa, kasih sayang, dukungan, dan pengorbanan yang senantiasa diberikan.
    \item Lembaga Pengelola Dana Pendidikan (LPDP), Kementerian Keuangan Republik Indonesia, atas dukungan pembiayaan selama penulis menempuh studi magister.
    \item Ibu Robithoh Annur dan Bapak Hari Purnama, selaku dosen pembimbing, atas bimbingan, arahan, kritik, dan masukan dalam penyempurnaan penelitian ini.
    \item Kekasih hati, Fitriani Luluk Ashari, atas dukungan, pengertian, semangat, dan pendampingan selama penyelesaian studi dan tesis ini.
    \item Rekan-rekan IOTA Research Group, atas diskusi, bantuan, dan kebersamaan selama proses penelitian.
    \item Rekan-rekan Magister Informatika Angkatan 2024 Ganjil, atas dukungan dan kebersamaan selama menjalani perkuliahan dan penyelesaian studi.
\end{enumerate}

Penulis menyadari bahwa tesis ini masih memiliki keterbatasan. 
Kritik dan saran yang membangun diharapkan dapat menjadi masukan bagi pengembangan penelitian selanjutnya. 
Semoga tesis ini bermanfaat bagi pengembangan ilmu pengetahuan, khususnya pada bidang keamanan siber, \textit{Attribute-Based Access Control}, dan komputasi \textit{edge-IoT}.

\vspace{2\baselineskip}

\begin{flushright}
Bandung, 18 September 2026

\vspace{2\baselineskip}

Penulis
\end{flushright}